\documentclass[10pt,a4paper]{amsart}
\usepackage[margin=19mm]{geometry}
\usepackage{amsmath,amssymb,mathtools}
\usepackage[scr=rsfs,cal=euler]{mathalfa}
\usepackage{xfrac}
\usepackage[unicode,hidelinks,bookmarks=false]{hyperref}
\newcommand{\ie}{i.e.\ }
\newcommand{\cf}{cf.\ }
\newcommand{\resp}{respectively\ }
\newcommand{\Subsection}{\subsection}
\providecommand{\nobreakdash}{\nobreak\hbox{-}}
\setlength{\emergencystretch}{2em}
\multlinegap0pt
\usepackage{xcolor}
\usepackage{tikz}
\usepackage{amssymb}
\usepackage{mathtools}
\usepackage{bm}
\newtheorem{lemma}{Lemma}
\def\eqref#1{equation~\ref{#1}}
\def\mS{{\bm{S}}}
\def\mSigma{{\bm{\Sigma}}}
\def\mW{{\bm{W}}}
\def\vu{{\bm{u}}}
\def\vzero{{\bm{0}}}
\title{How Does Attention Help? Insights from Random Matrices on Signal Recovery from Sequence Models}
\author{Mohamed El Amine Seddik}
\date{Source: arXiv:2605.06826v1 (CC BY 4.0)}
\begin{document}
\maketitle

\noindent\emph{Source and licence.} How Does Attention Help? Insights from Random Matrices on Signal Recovery from Sequence Models. Version arXiv:2605.06826v1 (CC BY 4.0), distributed by arXiv under the Creative Commons Attribution 4.0 International licence, http://creativecommons.org/licenses/by/4.0/, as recorded in the arXiv licence metadata for this exact version and shown on the abstract page https://arxiv.org/abs/2605.06826v1. Full licence notice: \texttt{licenses/arxiv-2605.06826-CC-BY-4.0.txt}.\par\smallskip

\noindent\textit{Editorial scope.} Counted unit: the article's lemma lem:overlap\_derivative (source0.tex lines 1315--1321). The article's complete original proof of it is delivered with it (source0.tex lines 1322--1346, one \texttt{proof} environment of 25 lines). No mathematics is written, repaired or re-derived: the statement and the proof are the article's own lines, in the article's own notation, and no retained line is substituted or re-flowed.  Retained uncounted as the source-local context the proof needs: lines 1309--1312.  Nothing was omitted from any retained span, so the package carries no omission record.  No retained line contains a citation, so the wrapper carries no bibliography and no bibliography entry.  No retained line contains a figure environment, an image-inclusion command or drawing code, so no image is delivered and the package declares no asset dependency.  Editorial formatting only, in addition: an AMS \texttt{amsart} preamble that restates the article's own declarations and macros used by the retained lines, namely the environments \texttt{lemma} on separate counters, together with the standard packages those lines need.  The retained environments print in this document as Lemma 1 in order of appearance, whereas the article prints the same items with the numbering of its own full document.  The complete native source of the article as distributed in the arXiv e-print for this exact version is bundled unchanged as \texttt{sources/200049/source0.tex}.
\begin{equation}\label{eq:Sigma_beta_decomp}
    \mSigma(\beta) \;=\; \beta\,\hat\vu_\mSigma\hat\vu_\mSigma^\top \;+\; \mSigma_\perp,
    \qquad \hat\vu_\mSigma^\top\mSigma_\perp=\vzero, \qquad \|\hat\vu_\mSigma\|=1.
\end{equation}

\begin{lemma}[Overlap/derivative identity]\label{lem:overlap_derivative}
    Under~\eqref{eq:Sigma_beta_decomp},
    \begin{equation}\label{eq:overlap_derivative}
        |\hat\vu_\mS(\beta)^\top\,\hat\vu_\mSigma|^2
        \;=\; \frac{\beta}{\lambda(\beta)}\,\frac{d\lambda(\beta)}{d\beta}.
    \end{equation}
\end{lemma}

\begin{proof}
In the eigenbasis of~\eqref{eq:Sigma_beta_decomp}, $\mSigma(\beta)^{1/2}=\sqrt\beta\,\hat\vu_\mSigma\hat\vu_\mSigma^\top + \mSigma_\perp^{1/2}$ and so
\[
    \frac{d}{d\beta}\mSigma(\beta)^{1/2} \;=\; \frac{1}{2\sqrt\beta}\,\hat\vu_\mSigma\hat\vu_\mSigma^\top.
\]
Differentiating $\mS(\beta)=\mSigma^{1/2}\mW\mSigma^{1/2}$,
\[
    \frac{d\mS}{d\beta} \;=\; \frac{d\mSigma^{1/2}}{d\beta}\mW\mSigma^{1/2} + \mSigma^{1/2}\mW\frac{d\mSigma^{1/2}}{d\beta}.
\]
For a simple eigenpair, standard first-order perturbation gives $d\lambda/d\beta = \hat\vu_\mS^\top (d\mS/d\beta)\hat\vu_\mS$. Substituting and using symmetry of $\mW$:
\[
    \frac{d\lambda}{d\beta} \;=\; \frac{1}{\sqrt\beta}\,\big(\hat\vu_\mSigma^\top\hat\vu_\mS\big)\,\hat\vu_\mSigma^\top\mW\mSigma^{1/2}\hat\vu_\mS.
\]
Now left-multiply the eigenequation $\mS\hat\vu_\mS=\lambda\hat\vu_\mS$ by $\hat\vu_\mSigma^\top$ and use $\hat\vu_\mSigma^\top\mSigma^{1/2}=\sqrt\beta\,\hat\vu_\mSigma^\top$:
\[
    \sqrt\beta\,\hat\vu_\mSigma^\top\mW\mSigma^{1/2}\hat\vu_\mS \;=\; \lambda\,\hat\vu_\mSigma^\top\hat\vu_\mS,
    \qquad\text{i.e.}\qquad
    \hat\vu_\mSigma^\top\mW\mSigma^{1/2}\hat\vu_\mS \;=\; \frac{\lambda}{\sqrt\beta}\,\hat\vu_\mSigma^\top\hat\vu_\mS.
\]
Plugging back,
\[
    \frac{d\lambda}{d\beta} \;=\; \frac{\lambda}{\beta}\,\big(\hat\vu_\mSigma^\top\hat\vu_\mS\big)^2,
\]
which rearranges to~\eqref{eq:overlap_derivative}.
\end{proof}

\end{document}
